\section*{Chapitre \ref{ch:g12:glucose-and-diabetes} --- La glycémie et le diabète}

\begin{solution}{exo:g12:glucose-and-diabetes:1}
Environ \qty{1}{g/L}, soit \qty{5.5}{mmol/L}. Au-dessous de
\qty{0.5}{g/L}, le cerveau défaille~; au-dessus de \qty{1.8}{g/L}, le
glucose passe dans l'urine (et des années au-dessus de \qty{1.3}{g/L}
abîment les vaisseaux).
\end{solution}

\begin{solution}{exo:g12:glucose-and-diabetes:2}
L'\omterm{prop:g12:glucose-and-diabetes:hormones}{insuline}, produite par les \omterm{def:g10:cells-common-unit:cell}{cellules} $\beta$, abaisse la \omterm{def:g12:glucose-and-diabetes:glycaemia}{glycémie}
(captage et mise en réserve par le muscle, la graisse et le foie)~; le
\omterm{prop:g12:glucose-and-diabetes:hormones}{glucagon}, produit par les \omterm{def:g10:cells-common-unit:cell}{cellules} $\alpha$, l'élève (libération de
glucose par le foie).
\end{solution}

\begin{solution}{exo:g12:glucose-and-diabetes:3}
Le foie, sous forme de \omterm{def:g10:exercise-and-energy:glycogen}{glycogène} (et il sait aussi fabriquer du
glucose neuf à partir d'\omterm{def:g10:chemistry-of-life:families}{acides aminés} et de lactate).
\end{solution}

\begin{solution}{exo:g12:glucose-and-diabetes:4}
Type 1~: les \omterm{def:g10:cells-common-unit:cell}{cellules} $\beta$ sont détruites et l'\omterm{prop:g12:glucose-and-diabetes:hormones}{insuline} est
absente. Type 2~: l'\omterm{prop:g12:glucose-and-diabetes:hormones}{insuline} est présente mais les tissus y répondent
mal, et les \omterm{def:g10:cells-common-unit:cell}{cellules} $\beta$ finissent par s'épuiser.
\end{solution}

\begin{solution}{exo:g12:glucose-and-diabetes:5}
Sujet sain environ \qty{0.95}{g/L}~; tolérance altérée environ
\qty{1.7}{g/L} (entre 1.4 et 2)~; type 2 environ \qty{2.7}{g/L}
(au-dessus de 2).
\end{solution}

\begin{solution}{exo:g12:glucose-and-diabetes:6}
Le foie libère du glucose à partir de son \omterm{def:g10:exercise-and-energy:glycogen}{glycogène} au rythme où le
cerveau le retire, sous le contrôle du \omterm{prop:g12:glucose-and-diabetes:hormones}{glucagon}, si bien que le
compartiment reste près de \qty{5}{g} bien que son contenu soit
renouvelé toutes les heures.
\end{solution}

\begin{solution}{exo:g12:glucose-and-diabetes:7}
Le canal conduit les \omterm{def:g11:enzymes-and-phenotype:enzyme}{enzymes} digestives à l'intestin~; l'obstruer
laisse la \omterm{def:g12:glucose-and-diabetes:glycaemia}{glycémie} normale, l'effet du pancréas sur le glucose ne
passe donc pas par la digestion. Retirer l'organe entier retire les
\omterm{prop:g12:glucose-and-diabetes:hormones}{îlots}~: l'\omterm{prop:g12:glucose-and-diabetes:hormones}{insuline} est fabriquée dans le pancréas et atteint ses
cibles par le sang, non par le canal.
\end{solution}

\begin{solution}{exo:g12:glucose-and-diabetes:8}
Le \omterm{prop:g12:glucose-and-diabetes:hormones}{glucagon} agit en faisant dégrader au foie son \omterm{def:g10:exercise-and-energy:glycogen}{glycogène}~; après une
nuit de jeûne, la réserve est là et la \omterm{def:g12:glucose-and-diabetes:glycaemia}{glycémie} monte~; après trois
jours, le \omterm{def:g10:exercise-and-energy:glycogen}{glycogène} est épuisé et le \omterm{prop:g12:glucose-and-diabetes:hormones}{glucagon} n'a plus rien à faire
libérer.
\end{solution}

\begin{solution}{exo:g12:glucose-and-diabetes:9}
Au-dessus de \qty{1.8}{g/L} environ, les reins ne peuvent pas
réabsorber tout le glucose qu'ils filtrent~; il passe dans l'urine et,
par osmose, y entraîne de l'eau. L'eau perdue doit être remplacée~:
c'est la soif.
\end{solution}

\begin{solution}{exo:g12:glucose-and-diabetes:10}
L'\omterm{prop:g12:glucose-and-diabetes:hormones}{insuline} pousse le glucose dans les muscles et les \omterm{def:g10:cells-common-unit:cell}{cellules}
adipeuses et arrête la libération par le foie, mais aucun glucose
n'arrive~: la \omterm{def:g12:glucose-and-diabetes:glycaemia}{glycémie} tombe au-dessous de \qty{0.6}{g/L} en une heure
ou deux, avec sueurs, tremblements et confusion~; du sucre pris
aussitôt corrige la chose.
\end{solution}

\begin{solution}{exo:g12:glucose-and-diabetes:11}
Un muscle qui travaille capte le glucose du sang par une voie qui n'a
pas besoin d'\omterm{prop:g12:glucose-and-diabetes:hormones}{insuline}, et un muscle entraîné regagne entre les séances
de la sensibilité à l'\omterm{prop:g12:glucose-and-diabetes:hormones}{insuline}~; la \omterm{def:g12:glucose-and-diabetes:glycaemia}{glycémie} baisse pendant et après
l'exercice, et la résistance des tissus, cause de la maladie, se
réduit.
\end{solution}

\begin{solution}{exo:g12:glucose-and-diabetes:12}
Pas d'\omterm{prop:g12:glucose-and-diabetes:hormones}{insuline}~: type 1, les \omterm{def:g10:cells-common-unit:cell}{cellules} $\beta$ des \omterm{prop:g12:glucose-and-diabetes:hormones}{îlots} sont détruites
et ne sécrètent rien. Trois fois l'\omterm{prop:g12:glucose-and-diabetes:hormones}{insuline} normale~: type 2, les
\omterm{prop:g12:glucose-and-diabetes:hormones}{îlots} sécrètent à plein régime contre des tissus qui ne répondent
plus.
\end{solution}

\begin{solution}{exo:g12:glucose-and-diabetes:13}
Sujet sain~: la \omterm{def:g12:glucose-and-diabetes:glycaemia}{glycémie} monte à environ \qty{1.3}{g/L} entre 30 et 60
minutes et revient en 2 heures~; l'\omterm{prop:g12:glucose-and-diabetes:hormones}{insuline} monte avec elle et
redescend avec elle. Type 1 non traité~: la \omterm{def:g12:glucose-and-diabetes:glycaemia}{glycémie} dépasse
\qty{2}{g/L} et y reste~; l'\omterm{prop:g12:glucose-and-diabetes:hormones}{insuline} est plate à zéro --- rien ne
retire le glucose. Type 2~: la \omterm{def:g12:glucose-and-diabetes:glycaemia}{glycémie} monte haut et redescend
lentement~; l'\omterm{prop:g12:glucose-and-diabetes:hormones}{insuline} monte plus haut que la normale et y
reste --- beaucoup d'\omterm{def:g11:hormones-and-reproduction:hormone}{hormone}, peu d'effet.
\end{solution}

\begin{solution}{exo:g12:glucose-and-diabetes:14}
$5 + 80 = \qty{85}{g}$ dans \qty{5}{L}~: \qty{17}{g/L}. Cela n'arrive
jamais parce que l'absorption prend une ou deux heures et que, à
mesure que le glucose arrive, l'\omterm{prop:g12:glucose-and-diabetes:hormones}{insuline} l'envoie dans le \omterm{def:g10:exercise-and-energy:glycogen}{glycogène} du
foie, celui des muscles et la graisse, tandis que le cerveau et les
tissus en brûlent une part~: le compartiment ne contient jamais plus
de quelques grammes au-dessus de la normale.
\end{solution}

\begin{solution}{exo:g12:glucose-and-diabetes:15}
Les deux ont un capteur (les \omterm{def:g10:cells-common-unit:cell}{cellules} des \omterm{prop:g12:glucose-and-diabetes:hormones}{îlots}~; l'hypothalamus et
l'hypophyse), une consigne, des effecteurs, et une \omterm{def:g11:hormones-and-reproduction:hormone}{rétroaction}
négative qui supprime le signal. La \omterm{def:g12:glucose-and-diabetes:glycaemia}{glycémie} dispose de deux \omterm{def:g11:hormones-and-reproduction:hormone}{hormones}
opposées agissant sur les mêmes organes, si bien que la grandeur peut
être poussée activement vers le haut comme vers le bas~; la boucle de
la \omterm{prop:g11:becoming-male-female:hormones}{testostérone} n'a qu'une \omterm{def:g11:hormones-and-reproduction:hormone}{hormone} et corrige surtout en en mettant
plus ou moins. Deux effecteurs opposés donnent un contrôle plus rapide
et plus serré d'une grandeur qui varie vite dans les deux sens.
\end{solution}

\begin{solution}{pb:g12:glucose-and-diabetes:1}
\textbf{1.} \qty{5}{g}~; $\qty{1}{g/L} \div \qty{180}{g/mol} = \qty{5.55}{mmol/L}$.

\textbf{2.} Une heure.

\textbf{3.} $100/5 = 20$ heures. Ensuite, le foie fabrique du glucose
neuf à partir d'\omterm{def:g10:chemistry-of-life:families}{acides aminés} (et de lactate), aux dépens des
\omterm{def:g10:chemistry-of-life:families}{protéines} du corps.

\textbf{4.} Les \omterm{def:g10:cells-common-unit:cell}{cellules} musculaires n'ont pas l'\omterm{def:g11:enzymes-and-phenotype:enzyme}{enzyme} qui libère du
glucose libre à partir du \omterm{def:g10:exercise-and-energy:glycogen}{glycogène} vers le sang~; leur \omterm{def:g10:exercise-and-energy:glycogen}{glycogène} est
brûlé sur place, pour leur propre contraction.

\textbf{5.} $5 + 60 = \qty{65}{g}$ dans \qty{5}{L}~: \qty{13}{g/L}.
L'\omterm{prop:g12:glucose-and-diabetes:hormones}{insuline} envoie le glucose dans le \omterm{def:g10:exercise-and-energy:glycogen}{glycogène} du foie et des muscles
et dans la graisse, et les tissus en brûlent une part~; la \omterm{def:g12:glucose-and-diabetes:glycaemia}{glycémie}
culmine autour de \qty{1.4}{g/L}.

\textbf{6.} À jeun \qty{1.15}{g/L} (au-dessous de 1.26) mais
\qty{1.85}{g/L} à 120 minutes (entre 1.4 et 2)~: tolérance altérée, au
seuil du \omterm{def:g12:glucose-and-diabetes:diabetes}{diabète}.

\textbf{7.} Non~: l'\omterm{prop:g12:glucose-and-diabetes:hormones}{insuline} est au double d'un pic normal. C'est la
réponse des tissus qui défaille --- une insulinorésistance.

\textbf{8.} Un \omterm{def:g12:glucose-and-diabetes:diabetes}{diabète} de type 2 (ou son stade d'alerte). À
attendre~: un adulte, en surpoids, sédentaire, souvent avec des
parents diabétiques.

\textbf{9.} Une \omterm{def:g12:glucose-and-diabetes:glycaemia}{glycémie} partant au-dessus de \qty{1.5}{g/L}, montant
au-dessus de \num{2.5} et ne redescendant pas~; une \omterm{prop:g12:glucose-and-diabetes:hormones}{insuline} proche de
zéro d'un bout à l'autre.

\textbf{10.} Une perte de poids et un exercice régulier~: ils
rétablissent la sensibilité des tissus à l'\omterm{prop:g12:glucose-and-diabetes:hormones}{insuline} et laissent le
muscle au travail capter du glucose sans elle.

\textbf{11.} $90/10 = 9$ unités.

\textbf{12.} La dose prend en charge \qty{90}{g} mais \qty{45}{g}
seulement arrivent~: l'\omterm{prop:g12:glucose-and-diabetes:hormones}{insuline} en excès fait descendre la \omterm{def:g12:glucose-and-diabetes:glycaemia}{glycémie} de
l'équivalent de \qty{45}{g} dans le compartiment et les
réserves --- bien au-dessous de \qty{0.6}{g/L}~: sueurs, tremblements,
confusion~; du sucre aussitôt.

\textbf{13.} $(5 + 90)/5 = \qty{19}{g/L}$ si rien n'était retiré~; en
réalité, plusieurs \unit{g/L}. Au-dessus de \qty{1.8}{g/L}, les reins
laissent passer le glucose dans l'urine, avec de l'eau~: urines
abondantes et soif.

\textbf{14.} Entre les repas, le foie libérerait sans cela du glucose
sans frein (le \omterm{prop:g12:glucose-and-diabetes:hormones}{glucagon} agit, rien ne s'y oppose) et les tissus
brûleraient de la graisse avec ses produits acides~: un fond
d'\omterm{prop:g12:glucose-and-diabetes:hormones}{insuline} est nécessaire à l'équilibre de base, et non aux seuls
repas.

\textbf{15.} Le lecteur mesure la \omterm{def:g12:glucose-and-diabetes:glycaemia}{glycémie} comme le faisait la \omterm{def:g10:cells-common-unit:cell}{cellule}
$\beta$~; le stylo délivre l'\omterm{prop:g12:glucose-and-diabetes:hormones}{insuline} qu'elle aurait sécrétée. Ils ne
peuvent pas remplacer l'ajustement continu, minute par minute~: la
patiente mesure quelques fois par jour et estime, là où la \omterm{def:g10:cells-common-unit:cell}{cellule}
mesurait sans cesse et répondait exactement.

\textbf{16.} Le type 2, dont les \omterm{def:g10:cells-common-unit:cell}{cellules} $\beta$ existent et peuvent
être poussées. Dans le type 1, il n'y a pas de \omterm{def:g10:cells-common-unit:cell}{cellule} $\beta$ à
stimuler.

\textbf{17.} La réponse des tissus à l'\omterm{prop:g12:glucose-and-diabetes:hormones}{insuline} --- le côté effecteur
de la boucle~: avec moins de graisse et des muscles actifs, un signal
d'\omterm{prop:g12:glucose-and-diabetes:hormones}{insuline} normal produit de nouveau un captage normal, et les
\omterm{def:g10:cells-common-unit:cell}{cellules} $\beta$ sont soulagées.

\textbf{18.} Chacune est sécrétée en réponse à l'écart opposé de la
même grandeur~: une \omterm{def:g12:glucose-and-diabetes:glycaemia}{glycémie} haute appelle l'\omterm{prop:g12:glucose-and-diabetes:hormones}{insuline} et fait taire le
\omterm{prop:g12:glucose-and-diabetes:hormones}{glucagon}, une \omterm{def:g12:glucose-and-diabetes:glycaemia}{glycémie} basse fait l'inverse. À \qty{1}{g/L}, l'îlot
sécrète les deux à un rythme faible et équilibré --- son état de
repos.

\textbf{19.} Une \omterm{def:g12:glucose-and-diabetes:glycaemia}{glycémie} basse tue le cerveau en quelques minutes,
une \omterm{def:g12:glucose-and-diabetes:glycaemia}{glycémie} haute abîme lentement sur des années~: le danger aigu a
plusieurs gardes indépendantes, le danger lent une
seule --- ce qui explique que la défaillance de la seule \omterm{prop:g12:glucose-and-diabetes:hormones}{insuline}
suffise à créer une maladie, et que le \omterm{def:g12:glucose-and-diabetes:diabetes}{diabète} soit fréquent.

\textbf{20.} Environ \qty{5}{g} de glucose dans le sang~; le \omterm{def:g10:exercise-and-energy:glycogen}{glycogène}
du foie couvre environ 20 heures~; le patient de la partie II a une
tolérance altérée, au bord d'un \omterm{def:g12:glucose-and-diabetes:diabetes}{diabète} de type 2~; la dose de la
partie III est de 9 unités.
\end{solution}
